% Created 2026-06-12 Fri 23:24
% Intended LaTeX compiler: pdflatex
\documentclass{article}
\usepackage[utf8]{inputenc}
\usepackage[T1]{fontenc}
\usepackage{graphicx}
\usepackage{longtable}
\usepackage{wrapfig}
\usepackage{rotating}
\usepackage[normalem]{ulem}
\usepackage{amsmath}
\usepackage{amssymb}
\usepackage{capt-of}
\usepackage{hyperref}

\usepackage{fancyhdr}
\usepackage[top=25mm, left=25mm, right=25mm]{geometry}
\usepackage{longtable}
\fancyhead[R]{}
\setlength\headheight{43.0pt}
\usepackage{tabularx}
\usepackage[T1]{fontenc}
\usepackage[utf8]{inputenc}
\usepackage[spanish]{babel}
\usepackage[backend=biber]{biblatex}
\addbibresource{/home/emilioalbanfs/G1ArqComp/G1ArqComp/bibliography/bibliography.bib}
\author{[Emilio Alban, Ariel Montufar, Karen Suarez, Samuel Churaco]}
\date{2026-06-12}
\title{Taller — Máquina IAS / Von Neumann\\\medskip
\large ICCD332 Arquitectura de Computadores}
\hypersetup{
 pdfauthor={[Emilio Alban, Ariel Montufar, Karen Suarez, Samuel Churaco]},
 pdftitle={Taller — Máquina IAS / Von Neumann},
 pdfkeywords={},
 pdfsubject={},
 pdfcreator={Emacs 27.1 (Org mode 9.7.5)}, 
 pdflang={Spanish}}
\begin{document}

\maketitle
\fancyhead[C]{\includegraphics[scale=0.05]{../images/logoEPN.jpg}\\
ESCUELA POLITÉCNICA NACIONAL\\FACULTAD DE INGENIERÍA DE SISTEMAS\\
ARQUITECTURA DE COMPUTADORES}
\thispagestyle{fancy}

\section{Objetivos}
\label{sec:orgcaa456f}

\begin{itemize}
\item Comprender el ciclo de captación, decodificación y ejecución de la
máquina IAS (arquitectura Von Neumann).
\item Identificar y aplicar el conjunto de instrucciones IAS para traducir
código máquina a lenguaje ensamblador y a notación RTL.
\item Simular mentalmente la ejecución de un programa en la máquina IAS
para determinar el estado final de los registros.
\end{itemize}

\section{Instrucciones}
\label{sec:orgb8e98bc}

Complete las solicitudes del taller y suba el archivo \texttt{.org} y \texttt{.pdf}
al aula virtual. Use las combinaciones de teclado para trabajar en Emacs.

\section{Recursos}
\label{sec:org8919834}

Para este taller consulte el libro \citetitle{stallings2006esp}
disponible en \href{https://epnecuador-my.sharepoint.com/:f:/g/personal/lenin\_falconi\_epn\_edu\_ec/EgjH2RoedD5NuswqpOt8ExsB\_DE052v9Rlrg0QpEtbimDg?e=WoqexR}{EPN-SHAREPOINT (Clic aquí)} y use su WSL con Emacs para editar el archivo.

\subsection{Observaciones y Recomendaciones}
\label{sec:org1e36609}
\begin{itemize}
\item Tenga en cuenta la ruta a las imágenes y a la bibliografía. Se le
recomienda que realice una clonación o un fork del repositorio de la
clase a fin de que no tenga problemas generando el documento.
\item Para la generación correcta de archivos la configuración básica debe
estar operativa y funcional (es decir, Emacs, \LaTeX).
\end{itemize}

\section{Taller: Máquina de Von Neumann}
\label{sec:orgadbe4c3}

La máquina de Von Neumann realiza operaciones en un ciclo repetitivo de
\textbf{captación}, \textbf{decodificación} y \textbf{ejecución}. Para esto, el computador
IAS usa una memoria de 40 bits. Cuando la memoria se usa
numéricamente, el dato representa un número en complemento a 2; donde
el bit inicial corresponde al signo. Cuando la memoria se usa para el
registro de instrucciones, se divide en dos partes de 20 bits cada una. Los 8
primeros bits de cada parte corresponden al \emph{opcode} y los restantes
12 bits de cada parte a la \emph{dirección (operando)} desde la que se debe
leer o a la que se debe escribir.

\subsection{Conjunto de Instrucciones IAS}
\label{sec:org2f1f4e1}

Complete la Tabla \ref{tab:tablaISA-IAS}  de instrucciones del computador IAS
\autocite[p.47]{stallings2006esp} usando la notación RTL\footnote{Register Transfer Language o Lenguaje de Transferencia de Registros.}. Se suministra
la Tabla en código \LaTeX{} para tal efecto.

Para editar la Tabla \ref{tab:tablaISA-IAS} ubíquese en
el bloque \texttt{\#+begin\_export latex} y ejecute el comando \texttt{C-c'}. Esto
dividirá el buffer en dos ventanas. El cursor salta a la nueva ventana
y le permite editar la tabla. Dentro de la nueva tabla el modo
predominante cambiará automáticamente a \LaTeX, lo que le permite
hacer cambios con la asistencia del IDE. Para salir y aceptar los
cambios ejecute de nuevo \texttt{C-c'}.

\begin{itemize}
\item \texttt{M} equivale a \emph{Alt}
\item \texttt{x} es la tecla \emph{x}
\item \texttt{RET} es presionar \emph{Enter}
\end{itemize}

Una vez activado el modo \LaTeX{} el siguiente código recibe colores
sobre las palabras clave.

\begin{table}[htbp]
  \caption{Instrucciones Máquina IAS}
  \label{tab:tablaISA-IAS}
  \centering
  \begin{tabular}{|cccc|}
    \hline
    Opcode   & Opcode Hex &  Símbolo           & RTL \\ \hline
    00001010 & 0x0A       &  LOAD MQ           & $[AC]\leftarrow[MQ]$\\
    00001001 & 0x09       &  LOAD MQ, M(X)     & $[MQ]\leftarrow[X]$\\
    00100001 & 0x21       &  STOR M(X)         & $[X]\leftarrow[AC]$\\
    00000001 & 0x01       &  LOAD M(X)         & $[AC]\leftarrow[X]$\\
    00000010 & 0x02       &  LOAD -M(X)        & $[AC]\leftarrow-[X]$\\        
    00000011 & 0x03       &  LOAD |M(X)|       & $[AC]\leftarrow|[X]|$\\
    00000100 & 0x04       &  LOAD -|M(X)|      & $[AC]\leftarrow-|[X]|$\\
    00001101 & 0x0D       &  JUMP M(X,0:19)    & $[PC]\leftarrowX(0:19)$ (IZQ)\\
    00001110 & 0x0E       &  JUMP M(X,20:39)   & $[PC]\leftarrowX(20:39$ (DER) \\
    00001111 & 0x0F       &  JUMP + M(X,0:19)  & $If[AC]\geq0,then[PC]\leftarrowX(0:19)$ (IZQ)\\
    00010000 & 0x10       &  JUMP + M(X,20:39) & $If[AC]\geq0,then[PC]\leftarrowX(20:39)$ (DER)\\
    00000101 & 0x05       &  ADD M(X)          & $[AC]\leftarrow[AC]+[X]$ \\
    00000111 & 0x07       &  ADD |M(X)|        & $[AC]\leftarrow[AC]+|[X]|$ \\
    00000110 & 0x06       &  SUB M(X)          & $[AC]\leftarrow[AC]-[X]$ \\
    00001000 & 0x08       &  SUB |M(X)|        & $[AC]\leftarrow[AC]-|[X]|$ \\
    00001011 & 0x0B       &  MUL M(X)          & $[AC,MQ]\leftarrow[MQ]\times[X]$ \\
    00001100 & 0x0C       &  DIV M(X)          & $[MQ,AC]\leftarrow[AC]\div[X]$ \\
    00010100 & 0x14       &  LSH               & $[AC]\leftarrow[AC]\times2$ \\
    00010101 & 0x15       &  RSH               & $[AC]\leftarrow[AC]\div2$ \\
    00010010 & 0x12       &  STOR M(X,8:19)    & $[X](8:19)\leftarrow[AC](28:39)$ \\
    00010011 & 0x13       &  STOR M(X,28:39)   & $[X](28:39)\leftarrow[AC](28:39)$ \\
    10011001 & 0x99       &  HALT              & $-----$ \\
    \hline
  \end{tabular}
\end{table}

\newpage

\subsection{Ejercicio}
\label{sec:orgea8d42f}

En la máquina IAS, las instrucciones se dividen en dos segmentos:
i) izquierdo, desde el bit 0 a 19, y ii) derecho, desde el bit 20 al 39. Primero
se ejecuta el lado izquierdo y luego el derecho. El
contador de programa inicia en la posición \([PC] \leftarrow 300\). El conjunto de
instrucciones del computador o ISA\footnote{Instruction Set Architecture: conjunto de instrucciones que la máquina IAS puede interpretar y ejecutar.} está definido en la Tabla
\ref{tab:tablaISA-IAS}. A fin de terminar el programa, se agrega la
instrucción \texttt{0x99} que lleva al computador al estado de HALT. Conteste las
siguientes preguntas considerando la Tabla \ref{tab:orgfed1a01} como el
mapa de memoria de la máquina y obtenga el programa en lenguaje ensamblador y en
RTL.

\begin{table}[htbp]
\caption{\label{tab:orgfed1a01}Mapa de Memoria}
\centering
\begin{tabular}{rrrrr}
\hline
Dirección & \(Opcode_1\) & \(X_1\) & \(Opcode_2\) & \(X_2\)\\[0pt]
\hline
0x300 & 01 & 940 & 06 & 941\\[0pt]
0x301 & 21 & 940 & 14 & 000\\[0pt]
0x302 & 21 & 941 & 02 & 940\\[0pt]
0x303 & 21 & 940 & 99 & 000\\[0pt]
\ldots{} & \ldots{} & \ldots{} & \ldots{} & \ldots{}\\[0pt]
0x940 & 00 & 000 & 00 & 005\\[0pt]
0x941 & 00 & 000 & 00 & 002\\[0pt]
\hline
\end{tabular}
\end{table}

\subsubsection{Preguntas}
\label{sec:org3852bfe}

\begin{enumerate}
\item ¿Qué resultado se tiene al final en el registro del acumulador?

Respuesta:
El valor final en el acumulador:
\begin{itemize}
\item En formato Decimal es -3
\item En formato Hexadecimal es FFFFFFFFFD
\([AC] \leftarrow FFFFFFFFFD\)
\end{itemize}

\item ¿Se sobrescribe algún registro como resultado de la ejecución? Si
es verdadero, indique qué registro y con qué valor.
Sí, tanto la memoria como el acumulador se sobrescriben durante la
ejecución.
Los valores finales en las direcciones de memoria modificadas son:
Respuesta:
\begin{itemize}
\item \([940] \leftarrow FFFFFFFFFD\)
\item \([941] \leftarrow 0000000006\)
\item \([AC] \leftarrow FFFFFFFFFD\)
\end{itemize}

\item ¿Cuál es el valor del bit de Carry? ¿Qué operación modifica el bit?
Respuesta:
\begin{itemize}
\item \(Carry = 0\)

\item Aunque después de la instrucción \(SUB M(941)\), es decir la
operación \(5-2\) generó un bit de carry \(= 1\), siguiente a esta
instrucción que es \(LSH\), es decir la operación \([AC]\times2\), el
bit de acarreo se actualizó a 0, ya que el bit anterior de carry
fue desplazado hacia afuera.
\end{itemize}
\end{enumerate}

\subsubsection{Código RTL}
\label{sec:org6361505}

Escriba en notación de transferencia de registros el programa que
ejecuta el computador IAS.
\begin{itemize}
\item \([AC]\leftarrow[940]\)
\item \([AC]\leftarrow[AC]-[941]\)
\item \([940]\leftarrow[AC]\)
\item \([AC]\leftarrow[AC]\times2\)
\item \([941]\leftarrow[AC]\)
\item \([AC]\leftarrow-[940]\)
\item \([940]\leftarrow[AC]\)
\item \(HALT\)
\end{itemize}


\subsubsection{Código Assembler}
\label{sec:org83af939}

Escriba el código Assembler del programa que ejecuta el computador IAS.

\begin{verbatim}
LOAD M(940)
SUB M(941)
STOR M(940)
LSH
STOR M(941)
LOAD-M(940)
STOR M(940)
HALT
\end{verbatim}

\section{Equipo de Trabajo}
\label{sec:org0587745}

Complete la información de los integrantes del grupo.

\begin{center}
\begin{tabular}{lll}
\hline
Nombre & Email & Rol\\[0pt]
\hline
Emilio Alban & emilio.alban@epn.edu.ec & Líder\\[0pt]
Ariel Montufar & ariel.montufar@epn.edu.ec & Integrante\\[0pt]
Karen Suarez & karen.suarez01@epn.edu.ec & Integrante\\[0pt]
Samuel Churaco & samuel.churaco@epn.edu.ec & Integrante\\[0pt]
\hline
\end{tabular}
\end{center}

\section{Verificación de Entregables [100\%]}
\label{sec:org73862db}

Ejecute \texttt{C-c C-c} sobre los ítems de tarea según se hayan cumplido o
no. Si un ítem no pudo realizarse, anote en la sección de problemas
las razones.

\begin{itemize}
\item[{$\boxtimes$}] Tabla de instrucciones IAS completada.
\item[{$\boxtimes$}] Preguntas del ejercicio respondidas.
\item[{$\boxtimes$}] Código RTL completo.
\item[{$\boxtimes$}] Código Assembler completo.
\item[{$\boxtimes$}] Revisión de ortografía con \texttt{ispell-buffer}.
\item[{$\boxtimes$}] Exportación a PDF con \texttt{M-x org-latex-export-to-pdf}.
\item[{$\boxtimes$}] Subir al aula virtual archivo \texttt{.org} y \texttt{.pdf}.
\end{itemize}

\subsection{Problemas con la tarea / logísticos}
\label{sec:orgd3e934d}

\begin{itemize}
\item Tarea \(X_1\) no pudo completarse debido a: N/A
\item Tarea \(X_2\) no pudo completarse debido a: N/A
\end{itemize}

\printbibliography
\end{document}
